\chapter{Carbon capture}

\section{Introduction: Context and Motivation}

Atmospheric CO\textsubscript{2} concentrations have now permanently exceeded 400~ppm, 
accompanied by a measurable rise in global temperatures. This situation calls for a 
combination of mitigation strategies, including reductions in energy consumption, 
substitution of fossil fuels with natural gas or biomass, deployment of renewable and 
nuclear energy, and Carbon Capture, Utilization, and Storage (CCUS). Among these, CCUS 
stands out as one of the few approaches capable of addressing emissions from 
hard-to-decarbonize industrial sectors at a large scale.

Industrial flue gases represent the primary target for capture technologies. Depending 
on the source, the CO\textsubscript{2} molar fraction (y\textsubscript{CO\textsubscript{2}}) 
ranges from approximately 3\% in combined-cycle natural gas plants (CCGT) to around 
15\% in coal-fired power stations. Lime and cement plants are particularly CO\textsubscript{2}-intensive 
due to the calcination reaction (CaCO\textsubscript{3} $\rightarrow$ CaO + CO\textsubscript{2}), 
which can drive y\textsubscript{CO\textsubscript{2}} up to 20--30\%, making them prime 
candidates for capture retrofits.

A complementary approach is \textbf{Direct Air Capture (DAC)}, which extracts 
CO\textsubscript{2} directly from ambient air (at $\approx$400~ppm). While energetically 
more demanding---requiring 19--21~kJ/mol CO\textsubscript{2} compared to 1--9~kJ/mol 
for industrial flue gas sources---DAC offers a key land-use advantage: it requires 
500 to 1500 times less surface area than forests to capture an equivalent quantity of 
CO\textsubscript{2}. A well-known example is the ORCA facility in Iceland, capturing 
4000~t\textsubscript{CO\textsubscript{2}}/year using renewable geothermal energy. In 
the CCUS value chain, captured CO\textsubscript{2} is transported by pipeline or ship 
and either stored geologically (with an estimated worldwide capacity of 2000--20\,000 
Gt\textsubscript{CO\textsubscript{2}}) or utilized in a growing range of industrial, 
chemical, and food applications, with a utilization potential estimated at 
1--7~Gt\textsubscript{CO\textsubscript{2}}/year.

The \textbf{EU Emissions Trading Scheme (ETS)} provides an economic incentive: with 
CO\textsubscript{2} allowance prices currently above 70~\euro/t\textsubscript{CO\textsubscript{2}}, 
large emitters face a financial pressure that makes investment in capture technologies 
increasingly competitive.

\section{CO\texorpdfstring{\textsubscript{2}}{2} Capture Strategies}

Three broad combustion-side strategies can be distinguished:

\begin{itemize}
    \item \textbf{Pre-combustion capture:} The fuel (natural gas, coal, or biomass) is 
    first reformed or gasified, then subjected to a water-gas shift reaction, producing 
    a high-pressure H\textsubscript{2}/CO\textsubscript{2} mixture. The CO\textsubscript{2} 
    is separated before combustion, and the purified hydrogen is subsequently burned.
    
    \item \textbf{Oxyfuel combustion:} The fuel is burned in a nearly pure oxygen 
    atmosphere, yielding a flue gas composed predominantly of CO\textsubscript{2} and 
    water (>80\% CO\textsubscript{2}). Water is condensed, leaving a concentrated 
    CO\textsubscript{2} stream ready for compression and storage with minimal further 
    treatment.
    
    \item \textbf{Post-combustion capture:} CO\textsubscript{2} is removed from the 
    flue gas after a conventional combustion process. This so-called ``end-of-pipe'' 
    approach is the most technologically mature and the most easily retrofittable to 
    existing industrial facilities, as it does not require modification of the 
    combustion unit itself.
\end{itemize}

For each of these strategies, several physico-chemical separation techniques exist: 
gas-liquid absorption, gas-solid adsorption, membrane separation (gas-gas or 
gas-liquid), cryogenic condensation, calcium looping, chemical looping combustion, 
and biological approaches using micro-algae.

\section{Gas-Liquid Absorption: The Reference Technology}

\subsection{Process Description}

Gas-liquid (G-L) absorption remains the most industrially proven technology for 
post-combustion CO\textsubscript{2} capture. The process, originally patented by 
R.~R.~Bottoms in 1930 for acid gas separation, relies on two coupled columns:

\begin{itemize}
    \item An \textbf{absorber}, operating at approximately 40\textdegree C and 
    1.2~bar, where the flue gas contacts a descending lean solvent and CO\textsubscript{2} 
    is chemically or physically absorbed.
    \item A \textbf{regenerator (stripper)}, operating at around 120\textdegree C and 
    2~bar, where heat is applied to release concentrated CO\textsubscript{2} and 
    regenerate the solvent for recirculation.
\end{itemize}

The \textbf{regeneration energy} has three contributions:
\[
E_{\text{regen}} = \Delta H_{\text{Reaction}} + \Delta H_{\text{Sensible}} + \Delta H_{\text{VAP, solvent}}
\]
where $\Delta H_{\text{Reaction}}$ is the energy needed to break the solvent--CO\textsubscript{2} 
bond, $\Delta H_{\text{Sensible}}$ is the energy required to heat the solvent to the 
stripping temperature, and $\Delta H_{\text{VAP, solvent}}$ accounts for solvent 
vaporization. For a reference 30\% monoethanolamine (MEA) solution, this total is 
approximately 3.7~GJ/t\textsubscript{CO\textsubscript{2}}.

\subsection{Solvent Generations}

Solvent selection is governed by multiple interacting criteria: reaction kinetics with 
CO\textsubscript{2}, physico-chemical properties (density, viscosity, diffusivity, 
solubility), regeneration energy, solvent degradation behavior with CO\textsubscript{2}, 
oxygen, SO\textsubscript{x} and NO\textsubscript{x}, industrial availability, and cost.

\textbf{First-generation solvents} are single-component aqueous solutions, typically 
at 30\% concentration:
\begin{itemize}
    \item \textit{Primary amines} such as MEA react rapidly via a zwitterion intermediate 
    to form carbamates, offering high reactivity but high regeneration energy and 
    significant degradation susceptibility.
    \item \textit{Tertiary amines} such as MDEA (N-methyldiethanolamine) react via 
    bicarbonate formation; they are slower but require less regeneration energy.
    \item \textit{Sterically hindered amines} (SHA) such as AMP 
    (2-amino-2-methyl-1-propanol) form very unstable carbamates that readily hydrolyze 
    to bicarbonate, combining higher loading capacity with lower regeneration energy.
    \item \textit{Cyclic diamines} such as piperazine (PZ) act as absorption 
    activators due to their fast kinetics, often used as promoters in blends.
\end{itemize}

\textbf{Second-generation solvents} are blends of two or more first-generation 
solvents, exploiting complementary properties. The activation mechanism---for example, 
MDEA activated by PZ (commercialized by BASF as aMDEA in the early 1980s) or the 
more recent CESAR-1 formulation (AMP + PZ)---significantly reduces regeneration energy 
while maintaining adequate absorption rates. Hybrid solvents combining chemical and 
physical absorbents (e.g., Sulfinol) are also in this category.

\textbf{Third-generation solvents} include ionic liquids (ILs), which are molten salts 
liquid at ambient conditions with very low vapor pressure and high thermal stability 
(e.g., imidazolium, pyridinium, ammonium families), and amino acid salts such as 
potassium prolinate (KPro) or potassium sarcosinate (Ksar), which replace a hydrogen 
atom of classical amines by a cation, yielding non-volatile, potentially biodegradable 
solvents.

\textbf{Fourth-generation solvents} represent the current research frontier and include:
\begin{itemize}
    \item \textit{Enzymatic solvents:} A conventional solvent (amine or K\textsubscript{2}CO\textsubscript{3}) 
    is combined with the enzyme carbonic anhydrase, which catalyzes the CO\textsubscript{2} 
    hydration reaction and dramatically accelerates the absorption rate.
    \item \textit{Demixing (phase-change) solvents:} Upon CO\textsubscript{2} absorption, 
    the loaded solvent splits into two liquid phases. Only the CO\textsubscript{2}-rich 
    phase is sent to the regenerator, substantially reducing the volume of liquid 
    requiring heating and thus lowering the regeneration energy by 20--30\% compared 
    to MEA. The DMX\textsuperscript{\textregistered} process developed by IFP Energies 
    Nouvelles and the DEEA-MAPA system are leading examples.
\end{itemize}

A comparison of specific reboiler duties across solvent generations illustrates the 
progress achieved: from 3.7--4.5~GJ/t\textsubscript{CO\textsubscript{2}} for MEA 
(circa 1990--2005) down to approximately 2.2~GJ/t\textsubscript{CO\textsubscript{2}} 
for the best phase-change solvent systems demonstrated around 2015.

\subsection{Process Improvements and Degradation}

Beyond solvent selection, process configuration plays a major role. Configurations 
studied include intercooling of the absorber, rich-solvent vapor recompression (RVC), 
lean vapor compression (LVC), split-flow arrangements, staged stripper feeding, 
partial flash regeneration, and multi-column schemes. Simulation results show that 
combining MDEA/PZ with rich vapor compression and intercooling can reduce 
E\textsubscript{regen} by up to \textbf{35\%} relative to the MEA base case, reaching 
a minimum of approximately 2.19~GJ/t\textsubscript{CO\textsubscript{2}}.

Packing selection also influences performance: structured packings (Flexipac, Mellapak, 
BX) consistently outperform random packings such as Raschig or Pall rings in terms of 
reboiler duty and solvent flow rate.

Solvent degradation is a major operational challenge. Mechanisms include carbamate 
polymerization, oxazolidinone and imidazolidinone formation via reaction with 
CO\textsubscript{2}, thermal degradation at regeneration temperatures, and oxidative 
degradation by O\textsubscript{2}, SO\textsubscript{x}, and NO\textsubscript{x} 
present in flue gas. Degradation products cause corrosion of carbon steel equipment 
and foaming in the absorber column, both of which increase operating costs and 
maintenance requirements.

\section{Other Capture Technologies}

\subsection{Gas-Solid Adsorption}

Adsorption relies on the physical attraction between CO\textsubscript{2} molecules and 
active sites at a solid adsorbent surface. Regeneration is achieved by either 
\textbf{Temperature Swing Adsorption (TSA)}---heating the bed to desorb 
CO\textsubscript{2}---or \textbf{(Vacuum) Pressure Swing Adsorption ((V)PSA)}---reducing 
total pressure to lower the CO\textsubscript{2} partial pressure and release it. Common 
adsorbents include activated carbons (derived from coconut, coal, or petroleum coke), 
zeolites (microporous alkali silicoaluminates), and Metal-Organic Frameworks (MOFs), 
which offer highly tunable pore structures and selectivities. These materials are 
currently at TRL~6 and represent an active research front.

\subsection{Membranes}

Two membrane configurations are relevant to CO\textsubscript{2} capture:

\textbf{Gas-gas membranes} act as semi-permeable barriers exploiting differences in 
molecular permeability. The driving force is the pressure difference across the 
membrane; CO\textsubscript{2} preferentially permeates to the low-pressure side. 
Membrane types range from cheap polymeric membranes (limited by thermal stability 
above 400\textdegree C) to more robust but expensive zeolite, palladium, or ceramic 
membranes. Their main advantages are compactness (no regeneration step required, 
smaller footprint than packed columns) and low maintenance. However, they are most 
competitive for small-to-medium gas flow rates and/or high CO\textsubscript{2} 
concentrations; at large industrial scales and low concentrations, amine absorption 
remains more cost-effective. Hybrid processes combining membranes and amines are 
particularly attractive at intermediate scales.

\textbf{Gas-liquid membranes} combine membrane and solvent technologies: a microporous 
membrane provides the gas-liquid interface without mixing the phases, while selectivity 
is governed entirely by the choice of absorption liquid on the permeate side. The 
membrane's sole role is to maximize interfacial area.

\subsection{Cryogenic Processes}

CO\textsubscript{2} can be separated by cooling the flue gas until it condenses as a 
liquid. Often preceded by a PSA step to enrich the CO\textsubscript{2} stream, as in 
Air Liquide's CRYOCAP\textsuperscript{TM}~FG process, this approach is best suited 
to streams already rich in CO\textsubscript{2} (>60\%) and at elevated pressure, such 
as in ammonia production. Its main disadvantages are a high refrigeration energy 
demand and the need to remove water and other condensable impurities upstream. Its 
chief advantage is the production of a highly pure liquid CO\textsubscript{2} ready 
for pipeline transport.

\subsection{Calcium Looping and Chemical Looping Combustion}

\textbf{Calcium Looping (CaL)} uses CaO as a regenerable solid sorbent. At 
$\approx$650\textdegree C in a carbonator, CaO reacts with CO\textsubscript{2} to 
form CaCO\textsubscript{3}. The loaded sorbent is then transferred to a calciner 
operating at $\approx$950\textdegree C, where it is regenerated and releases a 
concentrated CO\textsubscript{2} stream. The purged CaO can feed a cement kiln, 
offering synergies with the cement industry. A related technology developed by 
LEILAC uses indirect heated calcination to directly capture process CO\textsubscript{2} 
from cement and lime production without additional chemicals.

\textbf{Chemical Looping Combustion (CLC)} is an alternative to oxyfuel combustion 
in which a metal oxide (e.g., Fe\textsubscript{2}O\textsubscript{3}, NiO, CuO, 
Mn\textsubscript{3}O\textsubscript{4}) circulates between two fluidized-bed reactors. 
In the fuel reactor, the oxide is reduced by the fuel, producing CO\textsubscript{2} 
and H\textsubscript{2}O; water is easily condensed to yield pure CO\textsubscript{2}. 
In the air reactor, the reduced metal is re-oxidized, releasing nitrogen-depleted air. 
This inherently separates CO\textsubscript{2} without the energy cost of an air 
separation unit.

\subsection{Biological Approaches: Micro-algae}

Micro-algae consume CO\textsubscript{2} and light via photosynthesis in photobioreactors 
(open or closed). After separation from the culture medium, the algal biomass can 
be processed into biofuels or used in an ethanol production cycle. The main limitation 
remains the large land area required and the energy demand for mixing and harvesting, 
though algae can simultaneously remove other pollutants such as NO\textsubscript{x}.

\section{Industrial Deployment and Outlook}

As of 2025--2026, the global CCUS project pipeline has grown substantially: the 
Global CCS Institute reports 734 facilities in total (a 17\% year-on-year increase), 
with 77 in operation, 47 under construction, and a combined capture capacity of 
513~Mt\textsubscript{CO\textsubscript{2}}/year across all development stages---more 
than tripled since 2017.

\medskip

Notable industrial references include:

\begin{itemize}
    \item \textbf{Boundary Dam (Canada):} The world's first commercial post-combustion 
    CCS unit on a coal-fired power plant (SaskPower, Unit~3), capturing up to 
    1~Mt\textsubscript{CO\textsubscript{2}}/year and simultaneously removing 100\% 
    of SO\textsubscript{2} emissions. Operational data confirmed a 90\% CO\textsubscript{2} 
    capture rate, equivalent to removing approximately 250\,000 cars from the road 
    annually.
    
    \item \textbf{Niederaussem (Germany):} A pilot plant attached to a 965~MW lignite 
    power plant, treating $\approx$1550~Nm\textsuperscript{3}/h of flue gas at 15\% 
    CO\textsubscript{2}, achieving $\approx$90\% capture efficiency. The pilot has 
    validated MEA, BASF's OASE\textsuperscript{\textregistered}~blue (20\% lower 
    specific energy than MEA), and CESAR-1 (AMP + PZ) at packing heights of 16~m.
    
    \item \textbf{Technology Centre Mongstad (TCM), Norway:} The world's largest 
    CO\textsubscript{2} capture test center, providing open-access infrastructure 
    for testing both amine and chilled ammonia processes on real flue gases from a 
    combined heat and power plant and a refinery.
    
    \item \textbf{Norcem Brevik Cement Plant (Norway):} A 400\,000~t/year capture 
    demonstration project exploiting waste heat from the kiln for solvent regeneration 
    (50\% capture rate), now fully constructed---a first for the cement sector.
    
    \item \textbf{Duiven Waste Incineration Plant (Netherlands):} 165~t\textsubscript{CO\textsubscript{2}}/day 
    captured by amine absorption, illustrating scalability to waste-to-energy facilities.
\end{itemize}

Despite these advances, the primary remaining barrier to widespread deployment is 
economic rather than technical. Capture costs must decrease further to align with 
carbon prices and avoid excessive energy penalties. The pathway forward involves 
developing advanced solvents with lower regeneration energies, optimizing process 
configurations, scaling up novel capture concepts (membranes, adsorption, CaL), and 
building the CO\textsubscript{2} transport and storage infrastructure required to 
connect emitters with geological reservoirs.

\medskip

As a closing note, CCUS is not a substitute for emissions reduction at source. 
Renewable energy deployment, energy efficiency improvements, and demand reduction 
remain essential. CCUS provides a complementary, necessary tool for sectors where 
direct decarbonization is technically or economically constrained---and for removing 
historical CO\textsubscript{2} already in the atmosphere through BECCS or DAC. 
In the words of the author: \textit{``The best captured CO\textsubscript{2} is the 
one that we have not emitted.''}

\section{Podcast -- Lionel Dubois, Université de Mons (capture et utilisation du CO\texorpdfstring{\textsubscript{2}}{2})}

\subsubsection{Le Carbon Removal X Prize d'Elon Musk}
Elon Musk a lancé le \textit{Carbon Removal X Prize}, un appel à projets ouvert à tous, dont l'objectif est de trouver une solution efficace pour capturer du CO\textsubscript{2} -- que ce soit dans l'air, dans les océans ou ailleurs -- et de le stocker sur le long terme. Le prix total s'élève à 100 millions de dollars, répartis en plusieurs sous-prix, dont un \textit{milestone prize} composé de 15 récompenses intermédiaires d'un million de dollars chacune. Lionel Dubois souligne que, au regard de l'enjeu planétaire, ce montant reste insuffisant, mais que l'initiative a le mérite de faire bouger les lignes et de stimuler la réflexion.

\subsubsection{L'équipe internationale YUMET}
Dans le cadre du X Prize, Lionel Dubois a été invité à rejoindre une équipe internationale appelée YUMET (\textit{Y-U-M-E-T}), coordonnée par Brittany Zimmerman, une personnalité ayant notamment collaboré avec la NASA. Cette équipe rassemble des dizaines d'experts du monde entier, allant de spécialistes du CO\textsubscript{2} à des profils très éloignés du domaine. Lionel y participe à des réunions techniques de haut niveau, parfois en pleine nuit en raison du décalage horaire avec les États-Unis. Il retient de cette expérience un enrichissement intellectuel considérable, indépendamment du résultat final du concours.

\subsubsection{Remerciements à Lionel pour son cadeau Climeworks}
En guise d'illustration concrète de la compensation carbone, Lionel Dubois a offert en direct à son interlocuteur un cadeau symbolique : la capture de 10~kg de CO\textsubscript{2} dans l'air via la plateforme \textit{Climeworks}, une entreprise spécialisée dans la capture directe de CO\textsubscript{2} atmosphérique, notamment connue pour son site islandais Orca. Ce geste correspond à la neutralisation d'environ 100~km parcourus en voiture (sur la base de 100~g de CO\textsubscript{2} par kilomètre). Le coût de cette opération était de 1~€ par kilogramme, soit 1\,000~€ par tonne. Lionel en profite pour rappeler que la compensation carbone ne doit pas se substituer aux efforts de réduction des émissions en amont, mais peut constituer une composante parmi d'autres dans une stratégie globale.

\subsubsection{La cimenterie : un secteur \textit{hard to abate}}
Lionel Dubois a emmené son interlocuteur dans une cimenterie située près de Mons afin d'illustrer un cas concret d'industrie où les émissions de CO\textsubscript{2} sont en grande partie inévitables. La fabrication du ciment nécessite du clinker, lui-même produit à partir de carbonate de calcium (CaCO\textsubscript{3}). Lors de l'étape de décarbonatation, ce carbonate se décompose en chaux vive (CaO) et en CO\textsubscript{2}, ce qui représente environ deux tiers des émissions totales de la cimenterie -- et ce, quel que soit le combustible utilisé. Une cimenterie de taille moyenne produit environ 3\,000 tonnes de clinker par jour, et celle visitée en produisait jusqu'à 5\,000 tonnes par jour, correspondant à une émission équivalente de CO\textsubscript{2}. La cimenterie de demain pourrait adopter deux grandes approches : la post-combustion (ajout d'une unité de capture en sortie de fumées) ou l'oxycombustion (combustion enrichie en oxygène générant des fumées très riches en CO\textsubscript{2}, facilitant leur récupération).

\subsubsection{La thèse de Lionel Dubois}
La thèse de doctorat de Lionel Dubois, soutenue en 2013, portait sur la capture du CO\textsubscript{2} appliquée aux gaz de cimenterie, par la voie de la post-combustion. L'un des objectifs centraux était de concevoir et de mettre en \oe{}uvre en laboratoire un micro-pilote reproduisant les conditions réelles des fumées d'une cimenterie, afin de tester et comparer différentes solutions de capture et d'optimiser le procédé. Ce dispositif est toujours opérationnel aujourd'hui. La thèse comprenait également des aspects de simulation et d'essais en laboratoire, mais c'est le micro-pilote qui reste l'élément le plus marquant aux yeux de Lionel.

\subsubsection{L'échelle TRL (\textit{Technology Readiness Level})}
Les technologies de capture du CO\textsubscript{2} sont évaluées selon une échelle de maturité technologique appelée TRL (\textit{Technology Readiness Level}), allant de 1 (concept théorique) à 9 (technologie disponible commercialement). Parmi ces technologies, la plus mature -- à TRL~9 -- est l'absorption gaz-liquide (HAB), qui utilise un solvant aminé pour capturer chimiquement le CO\textsubscript{2} contenu dans les fumées. D'autres technologies, comme l'adsorption gaz-solide (TSA, PSA), les membranes ou les procédés cryogéniques, se trouvent à des niveaux de maturité inférieurs mais font l'objet de recherches actives. Le choix technologique doit être multicritère et adapté aux caractéristiques propres de chaque site industriel.

\subsubsection{Article sur la capture directe du CO\texorpdfstring{\textsubscript{2}}{2} dans les camions}
Lionel Dubois a été invité à commenter un article portant sur la capture du CO\textsubscript{2} directement à l'échappement de véhicules, à l'aide de matériaux adsorbants régénérés en cours de route. Il reconnaît l'intérêt scientifique de la démarche, mais émet des réserves : cette approche donne l'image d'une continuation du statu quo (on continue à brûler du carburant fossile tout en capturant le CO\textsubscript{2}), elle augmente le poids des véhicules à contre-courant des efforts d'allègement, et elle est moins pertinente à l'échelle de la voiture individuelle. Il se montrerait plus favorable à une approche amont : produire des carburants de synthèse à partir de CO\textsubscript{2} capturé, de sorte que le bilan soit neutre dès la fabrication du carburant. En revanche, l'application sur les grands navires lui semble plus prometteuse, car les volumes en jeu justifient davantage l'installation d'unités de capture compactes à bord.

\subsubsection{Le prix du CO\texorpdfstring{\textsubscript{2}}{2}}
Le coût de capture du CO\textsubscript{2} varie fortement selon le type d'application :
\begin{itemize}
    \item \textbf{Procédés de fermentation} (CO\textsubscript{2} déjà concentré) : environ 20 à 30~€ par tonne.
    \item \textbf{Applications industrielles} (cimenteries, centrales, etc.) : entre 50 et 80~€ par tonne pour la seule étape de capture, auxquels s'ajoutent les coûts de compression et de transport.
    \item \textbf{Capture directe dans l'air} (\textit{Direct Air Capture}) : au-delà de 100~€ par tonne, avec des chiffres pouvant atteindre 200 à 300~€ selon les estimations, voire 1\,000~€ à petite échelle (comme dans l'exemple Climeworks).
\end{itemize}
À titre de comparaison, le prix du quota de CO\textsubscript{2} dans le système ETS européen était d'environ 80~€ par tonne au moment de l'enregistrement du podcast. Lionel souligne que les coûts sont en diminution progressive grâce aux améliorations technologiques -- l'énergie nécessaire à la régénération des solvants est passée de 3,5~GJ/tonne à environ 2~GJ/tonne -- et que cette tendance devrait se poursuivre à mesure que davantage d'installations seront déployées.

\subsubsection{Les normes de qualité du CO\texorpdfstring{\textsubscript{2}}{2}}
La qualité du CO\textsubscript{2} capturé n'est pas uniforme et dépend fortement du procédé utilisé et de l'application envisagée. Pour un stockage géologique, une pureté élevée n'est pas indispensable. En revanche, pour une application alimentaire, le CO\textsubscript{2} doit être quasi pur (99,9999~\%). Les impuretés typiquement rencontrées sont des oxydes de soufre (SO\textsubscript{x}), des oxydes d'azote (NO\textsubscript{x}), du sulfure d'hydrogène (H\textsubscript{2}S), ou encore de l'eau résiduelle. Dans le cas d'un transport par pipeline, deux raisons imposent de respecter des seuils de pureté : d'une part, des raisons physico-chimiques (les impuretés modifient la viscosité et le comportement du mélange gazeux), d'autre part, des risques de corrosion des canalisations (notamment en présence d'eau et de composés acides). Des normes de qualité devront donc être définies et respectées par tous les émetteurs souhaitant accéder aux infrastructures communes de transport du CO\textsubscript{2}.